\SourceLeaf{187}und nun fast überzeugt sei, dass der euklidische Satz ohne irgendein datum nicht zu beweisen sei. Auch halte er es für nicht unwahrscheinlich, dass unsere Geometrie nur ein Kapitel einer allgemeineren sei. Ich erzählte ihm daraufhin, Sie hätten vor einigen Jahren öffentlich geäußert, dass man in dieser Sache im Grunde seit Euklids Zeiten nicht weitergekommen sei. Ja, Sie hätten mir mehrmals gesagt, dass es Ihnen auch nach eingehender Beschäftigung mit diesem Gegenstand nicht gelungen sei, die Absurdität einer solchen Annahme zu beweisen. --- Als er mir das erbetene Buch schickte, lag der beiliegende Zettel darin. Kurz darauf (Ende Dezember) bat er mich mündlich, Ihnen diesen Zettel bei Gelegenheit beizulegen und Sie in seinem Namen zu bitten, ihm bei Gelegenheit Ihr Urteil über seine Ideen mitzuteilen.

Abgesehen von allem Übrigen hat das Buch selbst den Vorzug, dass es reichhaltige Literatur zum Gegenstand enthält, die er, wie er mir sagt, auch weiterhin gesammelt hat. .\ .\ .\ .\ .\}

\bigskip
\noindent\hrulefill
\bigskip

\Needspace{10\baselineskip}
\begin{center}[\emph{\textls{Beilage:} Notiz von} \textsc{Schweikart}. \emph{Marburg, Dezember 1818.}]\end{center}

\{Es gibt eine zweifache Geometrie, --- eine Geometrie im engeren Sinn --- die \textls{euklidische}, und eine astralische Größenlehre.

Die Dreiecke der letzteren haben die Besonderheit, dass die Summe der drei Winkel nicht zwei rechten Winkeln entspricht.

Setzt man \textls{dies} voraus, lässt sich mit aller Strenge beweisen,

a) dass die Summe der 3 Winkel im Dreieck \textls{kleiner} als 2 rechte Winkel ist,

b) dass diese Summe umso kleiner wird, je mehr Fläche das Dreieck umfasst,

c) dass die Höhe eines gleichschenkligen rechtwinkligen Dreiecks zwar immer weiter zunimmt, je mehr man die Schenkel verlängert, dass sie aber eine gewisse Linie, die. ich die \textls{Konstante} nenne, nicht überschreiten kann.

